\documentclass[../../../../main.tex]{subfiles}
\begin{document}

\begin{flushright}
\footnotesize\textit{【自動文字起こし・要確認】}
\end{flushright}

{\bf [解]}
対称性から，立体のうち $y\ge0$ の部分の体積 $V_1$，もとめる体積 $V_2$ として，
\begin{align*}
V_2=2V_1 \quad \cdots \text{①}
\end{align*}
である．接線は $y$ 軸平行でないので，$m$ を傾きとして
\begin{align*}
\ell: y=m(x-2a)
\end{align*}
とおける．これが楕円と接するので，$X=x/a,\ Y=y/b$ とおき直した座標で，
\begin{align*}
bY=m(aX-2a)
\end{align*}
と $X^2+Y^2=1$ が接する．したがって
\begin{align*}
\frac{|2am|}{\sqrt{(am)^2+b^2}}=1
\end{align*}
各辺0以上から2乗して，
\begin{align*}
4a^2m^2=a^2m^2+b^2 \quad\therefore\ m=\pm\frac{\sqrt3}{3}\frac{b}{a}
\end{align*}
図の接線の傾きは負だから，$a,b>0$ より複号負をとって，
\begin{align*}
\ell: y=-\frac{\sqrt3}{3}\frac{b}{a}(x-2a)
\end{align*}
である．この時接点 $P\left(\dfrac{a}{2},\dfrac{\sqrt3}{2}b\right)$ となる．したがって，
\begin{align*}
V_1 = (\text{円錐部分}) - (\text{楕円部分}) \quad \cdots \text{②}
\end{align*}
において，
\begin{align*}
(\text{円錐部分}) = \frac{1}{3}\pi(2a)^2\frac{2\sqrt3}{3}b-\frac{1}{3}\pi\left(\frac{1}{2}a\right)^2\left(\frac{2\sqrt3}{3}b-\frac{\sqrt3}{2}b\right)
= \frac{1}{3}\pi\left[\frac{8\sqrt3}{3}a^2b-\frac{\sqrt3}{24}a^2b\right] = \frac{7\sqrt3}{8}\pi a^2b
\end{align*}
\begin{align*}
(\text{楕円部分}) = \pi\int_0^{\frac{\sqrt3}{2}b}a^2\left(1-\frac{y^2}{b^2}\right)dy=a^2\pi\left[y-\frac{1}{3b^2}y^3\right]_0^{\frac{\sqrt3}{2}b}=\frac{3\sqrt3}{8}\pi a^2b
\end{align*}
を②に代入して
\begin{align*}
V_1=\frac{7-3}{8}\sqrt3\pi a^2b=\frac{\sqrt3}{2}\pi a^2b
\end{align*}
①から
\begin{align*}
V_2=\sqrt3\pi a^2b
\end{align*}

\begin{figure}[htb]
\begin{center}
\begin{tikzpicture}[scale=0.9, >=stealth]
  \draw[->] (-0.3,0) -- (4.5,0) node[right] {$x$};
  \draw[->] (0,-0.3) -- (0,3) node[above] {$y$};
  \draw[thick] (0,0) ellipse (1.5cm and 1.3cm);
  \node[below] at (1.5,0) {$a$};
  \node[left] at (0,1.3) {$b$};
  \fill (3,0) circle (1pt) node[below] {$2a$};
  \draw[thick] (0,1.13) -- (3,0);
  \fill (0.75,1.13) circle (1pt) node[above] {$P$};
  \begin{scope}
    \clip (0,0) rectangle (3,1.3);
    \fill[pattern=north east lines, pattern color=gray!50] (1.5,0) -- (3,0) -- (0.75,1.13) arc (37:0:1.5 and 1.3) -- cycle;
  \end{scope}
\end{tikzpicture}
\end{center}
\caption{回転体の体積$V_1$，$V_2$を求めるための図}
\end{figure}


\newpage

\end{document}
